\documentclass[11pt]{article}
\usepackage[a4paper,margin=22mm]{geometry}
\usepackage{fontspec}
\setmainfont{DejaVu Serif}
\setsansfont{DejaVu Sans}
\usepackage{microtype}
\usepackage{xcolor}
\usepackage{hyperref}
\usepackage{enumitem}
\usepackage{parskip}
\definecolor{Accent}{HTML}{971D32}
\definecolor{Muted}{HTML}{666666}
\hypersetup{colorlinks=true,urlcolor=Accent,linkcolor=Accent}
\setlist{leftmargin=1.6em,itemsep=0.3em,topsep=0.35em}
\setcounter{tocdepth}{2}
\renewcommand{\contentsname}{Contents}
\newcommand{\sectionrule}{\par\vspace{-0.5em}\noindent\textcolor{Accent}{\rule{\linewidth}{0.6pt}}\par}
\begin{document}

\begin{center}
{\sffamily\bfseries\color{Accent} TOEFL WORD OF THE DAY\par}
\vspace{8mm}
{\fontsize{42}{48}\selectfont\bfseries linger\par}
\vspace{2mm}
{\Large\sffamily\color{Accent} /ˈlɪŋɡər/\par}
\vspace{2mm}
{\small\sffamily Local audio phrase: linger\par}
{\small\sffamily Audio file: audios/linger.mp3\par}
\vspace{2mm}
{\large\sffamily\color{Muted} intransitive verb\par}
\vspace{5mm}
{\sffamily stay longer than expected \textbullet\ continue to exist or have an effect \textbullet\ spend extra time on something\par}
\vspace{6mm}
{\large Linger can mean remaining somewhere longer than expected, continuing to exist after something ends, or spending extra time considering or enjoying something.\par}
\vspace{6mm}
\begin{minipage}{0.84\textwidth}
\itshape “The effects of the economic crisis lingered for several years.”
\end{minipage}\par
\vspace{10mm}
\textbf{Date:} August 23rd 2026\quad\textbullet\quad \textbf{Level:} B1--mid-B2 TOEFL
\vfill
Created and compiled by \textbf{Amir Kooshky}\par
\vspace{2mm}
\href{https://t.me/KooshkyTOEFL}{Telegram Channel}\quad\textbullet\quad
\href{https://www.instagram.com/kooshkytoefl}{Instagram}\quad\textbullet\quad
\href{https://t.me/KooshkyTOEFL\_pv}{Personal Telegram}
\end{center}
\newpage
\tableofcontents
\newpage

\section{Meanings and usage}
\sectionrule
\subsection{Meaning 1: Stay longer than expected}
\textbf{Part of speech:} intransitive verb

\textbf{IPA:} /ˈlɪŋɡər/

\textbf{Pronunciation phrase:} linger

\textbf{Local audio file:} audios/linger.mp3

\textbf{Register:} neutral; common in both everyday and formal writing

\textbf{Definition:} to stay somewhere longer than necessary or expected, often because you do not want to leave

\textbf{In simpler words:} stay for longer than expected

\textbf{Typical pattern:} linger + in, near, by, outside, or after + place or event

\subsubsection{Natural examples}
\begin{enumerate}
  \item Several students lingered after class to ask the professor additional questions.
  \item Visitors often linger in the museum's central hall before moving to the next exhibit.
  \item She lingered by the door for a few moments before finally leaving.
  \item A few guests lingered after the conference had officially ended.
  \item The tourists lingered near the ruins to watch the sunset.
  \item He lingered outside the office, unsure whether to go inside.
  \item Some customers linger in cafés for hours while working on their laptops.
  \item The audience lingered in the theater lobby to discuss the performance.
\end{enumerate}

\subsubsection{Nuance and implication}
\textbf{Central nuance:} Linger suggests remaining after the normal or expected time to leave.

\textbf{Usual implication:} The person may stay because of interest, hesitation, enjoyment, curiosity, or reluctance to leave.

\textbf{Compared with depart:} Depart means leave a place. Linger means remain there for some additional time instead of leaving promptly.

\subsubsection{Grammar notes}
\textbf{How to use it:} This meaning normally does not take a direct object. The person who stays is the subject.

\textbf{What can follow it:} It is often followed by a word or phrase showing place or time, such as in the room, near the entrance, or after class.

\textbf{Common preposition:} Common patterns include linger in a place, linger near or by something, and linger after an event.

\textbf{Position in a sentence:} Adverbs such as briefly, momentarily, or longer can be used to describe how long someone lingers.

\subsubsection{Useful patterns}
\textbf{Pattern:} linger in + place

\textbf{Use:} Use this when someone remains inside a place longer than expected.

\textbf{Example:} Several visitors lingered in the gallery after the tour ended.

\textbf{Pattern:} linger near or by + place or object

\textbf{Use:} Use this when someone stays close to a particular location.

\textbf{Example:} The children lingered near the entrance waiting for their parents.

\textbf{Pattern:} linger after + event

\textbf{Use:} Use this when someone stays after an activity has officially ended.

\textbf{Example:} A few participants lingered after the lecture.

\textbf{Pattern:} linger outside + place

\textbf{Use:} Use this when someone waits or remains outside rather than entering or leaving.

\textbf{Example:} Reporters lingered outside the courthouse.

\subsubsection{Common wrong usage}
\paragraph{Correction 1}
\textbf{Wrong:} The students lingered the classroom after class.

\textbf{Correct:} The students lingered in the classroom after class.

\textbf{Why:} Linger is normally intransitive, so use a preposition such as in before the place.

\paragraph{Correction 2}
\textbf{Wrong:} She lingered to the door for several minutes.

\textbf{Correct:} She lingered by the door for several minutes.

\textbf{Why:} Use by or near for the place where someone remains.

\paragraph{Correction 3}
\textbf{Wrong:} The guests were lingered after dinner.

\textbf{Correct:} The guests lingered after dinner.

\textbf{Why:} Use linger directly in the active form. Do not add were.

\subsection{Meaning 2: Continue to exist or have an effect}
\textbf{Part of speech:} intransitive verb

\textbf{IPA:} /ˈlɪŋɡər/

\textbf{Pronunciation phrase:} linger

\textbf{Local audio file:} audios/linger.mp3

\textbf{Register:} neutral to formal; common in academic, historical, social, medical, and descriptive writing

\textbf{Definition:} to continue to exist, be noticeable, or have an effect after the original cause or event has ended

\textbf{In simpler words:} remain or continue after something ends

\textbf{Typical pattern:} linger + in, among, after, for, or on

\subsubsection{Natural examples}
\begin{enumerate}
  \item The smell of smoke lingered in the room long after the fire had been extinguished.
  \item Doubts about the safety of the product continued to linger among consumers.
  \item The effects of the economic crisis lingered for several years.
  \item A sense of uncertainty lingered after the meeting.
  \item The memory of the disaster lingered in the community for decades.
  \item Tension between the two groups continued to linger despite the agreement.
  \item The taste of the medicine lingered in his mouth.
  \item Questions about the reliability of the data still linger.
\end{enumerate}

\subsubsection{Nuance and implication}
\textbf{Central nuance:} This sense emphasizes persistence: something remains even though the event or cause that produced it is already over.

\textbf{Usual implication:} The remaining thing may be physical, such as a smell or taste, or abstract, such as doubt, fear, memory, tension, or an economic effect.

\textbf{Compared with disappear:} Something that disappears stops being present. Something that lingers remains noticeable or influential for some time.

\subsubsection{Grammar notes}
\textbf{How to use it:} The thing that continues to exist is the subject, so this use does not normally take a direct object.

\textbf{What can follow it:} Common subjects include smell, taste, memory, doubt, question, fear, tension, effect, and uncertainty.

\textbf{Common preposition:} Use linger in for a place or memory, linger among for a group of people, linger after for an event, and linger for for a length of time.

\textbf{Position in a sentence:} Still and continue to frequently appear with linger when emphasizing persistence.

\subsubsection{Useful patterns}
\textbf{Pattern:} a smell or taste lingers

\textbf{Use:} Use this when a physical sensation remains noticeable.

\textbf{Example:} The smell of paint lingered in the hallway.

\textbf{Pattern:} doubts or questions linger

\textbf{Use:} Use this when uncertainty remains unresolved.

\textbf{Example:} Questions linger about the long-term effects of the treatment.

\textbf{Pattern:} effects linger

\textbf{Use:} Use this when the consequences of an event continue after the event itself.

\textbf{Example:} The effects of the drought lingered for several seasons.

\textbf{Pattern:} linger for + period of time

\textbf{Use:} Use this to say how long something continues.

\textbf{Example:} The economic effects lingered for more than a decade.

\textbf{Pattern:} continue to linger

\textbf{Use:} Use this to strongly emphasize that something has still not disappeared.

\textbf{Example:} Concerns about data privacy continue to linger.

\subsubsection{Common wrong usage}
\paragraph{Correction 1}
\textbf{Wrong:} The smell lingered the room.

\textbf{Correct:} The smell lingered in the room.

\textbf{Why:} Use in before the place where a smell or effect remains.

\paragraph{Correction 2}
\textbf{Wrong:} The effects of the crisis lingered during five years.

\textbf{Correct:} The effects of the crisis lingered for five years.

\textbf{Why:} Use for to state a duration.

\paragraph{Correction 3}
\textbf{Wrong:} Many doubts are still lingering them.

\textbf{Correct:} Many doubts are still lingering.

\textbf{Why:} Linger is intransitive here and does not take a direct object.

\subsection{Meaning 3: Spend extra time on something}
\textbf{Part of speech:} intransitive verb

\textbf{IPA:} /ˈlɪŋɡər/

\textbf{Pronunciation phrase:} linger

\textbf{Local audio file:} audios/linger.mp3

\textbf{Register:} neutral to formal; common in descriptions of writing, discussion, film, reading, and leisurely activities

\textbf{Definition:} to spend more time than usual considering, looking at, or enjoying something

\textbf{In simpler words:} spend extra time looking at or thinking about something

\textbf{Typical pattern:} linger over or linger on + subject, detail, image, question, or activity

\subsubsection{Natural examples}
\begin{enumerate}
  \item The report lingers over the causes of the conflict but says little about its consequences.
  \item She lingered over her coffee while reading the morning newspaper.
  \item The documentary lingers on several images of the damaged landscape.
  \item The author lingers over small details that reveal the character's personality.
  \item He lingered over the final question because it required careful thought.
  \item The camera lingers on the empty street before the scene changes.
  \item The lecturer did not linger on the historical background because time was limited.
  \item Readers may want to linger over the final chapter because it contains several complex arguments.
\end{enumerate}

\subsubsection{Nuance and implication}
\textbf{Central nuance:} This use suggests giving something more time or attention than is strictly necessary.

\textbf{Usual implication:} The extra time may come from enjoyment, interest, careful thought, or deliberate emphasis.

\subsubsection{Grammar notes}
\textbf{How to use it:} Use over or on before the thing receiving extra time or attention.

\textbf{What can follow it:} Common objects after over or on include a meal, question, detail, topic, image, memory, or point.

\textbf{Common preposition:} Linger over often suggests spending extra time considering or enjoying something. Linger on often emphasizes continued attention to a topic, detail, or image.

\textbf{Position in a sentence:} The subject may be a person, writer, report, camera, or discussion.

\subsubsection{Useful patterns}
\textbf{Pattern:} linger over + meal, drink, question, or detail

\textbf{Use:} Use this when someone spends extra time enjoying or considering something.

\textbf{Example:} The students lingered over the final question before submitting their answers.

\textbf{Pattern:} linger on + topic, image, or point

\textbf{Use:} Use this when attention remains focused on something for longer than usual.

\textbf{Example:} The documentary lingers on images of abandoned factories.

\textbf{Pattern:} not linger on + topic

\textbf{Use:} Use this when someone chooses not to spend much time discussing something.

\textbf{Example:} The professor mentioned the disagreement but did not linger on it.

\subsubsection{Common wrong usage}
\paragraph{Correction 1}
\textbf{Wrong:} The report lingers about the causes of the crisis.

\textbf{Correct:} The report lingers over the causes of the crisis.

\textbf{Why:} Use linger over when extra time or attention is given to a subject.

\paragraph{Correction 2}
\textbf{Wrong:} The camera lingered the image for several seconds.

\textbf{Correct:} The camera lingered on the image for several seconds.

\textbf{Why:} Use linger on before the thing receiving continued attention.

\section{Word family}
\sectionrule
\subsection{lingering}
\textbf{Part of speech:} adjective

\textbf{IPA:} /ˈlɪŋɡərɪŋ/

\textbf{Definition:} continuing for longer than expected or remaining after the main event or cause has ended

\textbf{Relationship to the headword:} This adjective describes something that continues to remain or have an effect.

\textbf{Grammar pattern:} lingering + doubt, concern, effect, question, smell, memory, or problem

\textbf{Usage note:} This is especially common with abstract problems and effects that have not completely disappeared.

\subsubsection{Examples}
\begin{enumerate}
  \item The country continues to face the lingering effects of the recession.
  \item There are lingering doubts about the accuracy of the original study.
  \item A lingering smell of smoke remained in the building.
\end{enumerate}

\section{Common collocations}
\sectionrule
\subsection{linger after}
\textbf{Applies to:} Stay longer than expected

\textbf{Meaning:} remain after an event has ended

\textbf{Pattern:} linger after + event

\textbf{Example:} Several students lingered after the lecture to speak with the professor.

\subsection{linger near}
\textbf{Applies to:} Stay longer than expected

\textbf{Meaning:} remain close to a particular place

\textbf{Pattern:} linger near + place

\textbf{Example:} Visitors lingered near the entrance while waiting for the tour to begin.

\subsection{linger in}
\textbf{Applies to:} Continue to exist or have an effect

\textbf{Meaning:} remain present in a place

\textbf{Pattern:} linger in + place

\textbf{Example:} The smell of smoke lingered in the room for hours.

\subsection{doubts linger}
\textbf{Applies to:} Continue to exist or have an effect

\textbf{Meaning:} uncertainty continues to exist

\textbf{Example:} Doubts linger about whether the policy will achieve its goals.

\subsection{questions linger}
\textbf{Applies to:} Continue to exist or have an effect

\textbf{Meaning:} questions remain unresolved

\textbf{Example:} Important questions linger about the long-term environmental effects.

\subsection{effects linger}
\textbf{Applies to:} Continue to exist or have an effect

\textbf{Meaning:} consequences continue after their original cause

\textbf{Example:} The effects of prolonged unemployment can linger for years.

\subsection{memory lingers}
\textbf{Applies to:} Continue to exist or have an effect

\textbf{Meaning:} a memory remains strong for a long time

\textbf{Example:} The memory of the disaster still lingers in the community.

\subsection{linger for years}
\textbf{Applies to:} Continue to exist or have an effect

\textbf{Meaning:} continue for a long period

\textbf{Pattern:} linger for + period of time

\textbf{Example:} Political tensions created by the conflict lingered for years.

\subsection{lingering doubts}
\textbf{Applies to:} Continue to exist or have an effect

\textbf{Meaning:} doubts that continue to remain

\textbf{Example:} The new evidence removed some lingering doubts about the theory.

\textbf{Usage note:} Lingering is the adjective form.

\subsection{lingering effects}
\textbf{Applies to:} Continue to exist or have an effect

\textbf{Meaning:} effects that continue long after their original cause

\textbf{Example:} The report examines the lingering effects of childhood poverty.

\textbf{Usage note:} Lingering is the adjective form.

\subsection{linger over}
\textbf{Applies to:} Spend extra time on something

\textbf{Meaning:} spend extra time considering, discussing, or enjoying something

\textbf{Pattern:} linger over + noun

\textbf{Example:} The author lingers over the social consequences of industrialization.

\subsection{linger on}
\textbf{Applies to:} Spend extra time on something

\textbf{Meaning:} keep attention focused on something

\textbf{Pattern:} linger on + topic, detail, or image

\textbf{Example:} The camera lingers on the damaged building before moving to the next scene.

\section{Synonyms and antonyms}
\sectionrule
\subsection{Close synonyms}
\subsubsection{remain}
\textbf{Part of speech:} intransitive verb

\textbf{Applies to:} Stay longer than expected

\textbf{Shared meaning:} Both can mean continue to stay in a place.

\textbf{Difference:} Remain simply means continue to stay somewhere. Linger adds the idea of staying longer than expected, often because of hesitation, interest, or reluctance to leave.

\textbf{Example:} Several guests remained in the hall after the ceremony.

\subsubsection{persist}
\textbf{Part of speech:} intransitive verb

\textbf{Applies to:} Continue to exist or have an effect

\textbf{Shared meaning:} Both can describe something that continues to exist.

\textbf{Difference:} Persist emphasizes continuing despite time, difficulty, or attempts to stop something. Linger often suggests gradually remaining after the main cause or event has ended.

\textbf{Example:} The symptoms persisted despite treatment.

\subsubsection{endure}
\textbf{Part of speech:} intransitive verb

\textbf{Applies to:} Continue to exist or have an effect

\textbf{Shared meaning:} Both can describe something continuing over time.

\textbf{Difference:} Endure often suggests lasting for a long time and sometimes surviving difficult conditions. Linger more often describes a remaining effect, feeling, smell, memory, or doubt.

\textbf{Example:} The tradition endured for centuries.

\subsubsection{dwell on}
\textbf{Part of speech:} verb phrase

\textbf{Applies to:} Spend extra time on something

\textbf{Shared meaning:} Both can describe giving a subject extended attention.

\textbf{Difference:} Dwell on usually means think or talk about something for a long time, often something unpleasant. Linger on or over is broader and can involve thought, discussion, visual attention, or enjoyment.

\textbf{Example:} The report does not dwell on the causes of the failure.

\subsection{Direct or useful antonyms}
\subsubsection{depart}
\textbf{Part of speech:} intransitive verb

\textbf{Applies to:} Stay longer than expected

\textbf{Opposition:} Depart means leave a place, while linger means remain there for additional time.

\textbf{Example:} Most visitors departed immediately after the ceremony.

\subsubsection{disappear}
\textbf{Part of speech:} intransitive verb

\textbf{Applies to:} Continue to exist or have an effect

\textbf{Opposition:} Disappear means stop being present or noticeable, while linger means continue to remain noticeable.

\textbf{Example:} The smell gradually disappeared after the windows were opened.

\section{Words learners may confuse}
\sectionrule
\subsection{persist}
\textbf{Part of speech:} intransitive verb

\textbf{Why learners confuse them:} Both can mean continue to exist for longer than expected.

\textbf{Key difference:} Linger often describes something remaining after its original cause has ended, while persist more strongly emphasizes continuing despite time, difficulty, or attempts to stop it.

\textbf{linger:} A faint smell of smoke lingered in the room.

\textbf{persist:} The infection persisted despite several weeks of treatment.

\subsection{loiter}
\textbf{Part of speech:} intransitive verb

\textbf{Why learners confuse them:} Both can describe remaining in a place instead of leaving.

\textbf{Key difference:} Linger can simply mean remain somewhere because of interest, hesitation, or enjoyment. Loiter usually means stay in a public place without a clear purpose and can have a negative or legal implication.

\textbf{linger:} Visitors lingered in the gallery to admire the paintings.

\textbf{loiter:} The sign warned people not to loiter outside the entrance.

\section{Etymology}
\sectionrule
\textbf{Origin:} Linger comes from Middle English and has long carried the idea of remaining, delaying, or staying longer.

\textbf{Memory link:} If something lingers, it stays behind instead of disappearing or moving on quickly.

\section{Practice exercises}
\sectionrule
\subsection{Word family choice}
Choose the best member of the word family for each blank.

\begin{enumerate}
  \item The country is still dealing with the \_\_\_\_\_ effects of the financial crisis.
  \item Several important questions continue to \_\_\_\_\_ about the study's methodology.
\end{enumerate}

\subsection{Correct or incorrect?}
Decide whether each sentence is correct. Correct the inaccurate ones.

\begin{enumerate}
  \item The smell of smoke lingered in the room for several hours.
  \item The students lingered the classroom after the lesson.
  \item Doubts about the accuracy of the data continue to linger.
  \item The economic effects lingered during several years.
  \item The documentary lingers on images of the abandoned city.
\end{enumerate}

\subsection{Collocation practice}
Complete each sentence with a natural collocation.

\begin{enumerate}
  \item The effects of the recession continued to \_\_\_\_\_ for several years.
  \item A few students lingered \_\_\_\_\_ class to speak with the professor.
  \item Important doubts still \_\_\_\_\_ about the reliability of the results.
  \item The author lingers \_\_\_\_\_ the social consequences of the policy.
  \item The report examines the \_\_\_\_\_ effects of prolonged unemployment.
\end{enumerate}

\subsection{Complete answer key}
\subsubsection{Word family choice}
\begin{enumerate}
  \item \textbf{lingering.} The country is still dealing with the lingering effects of the financial crisis. \emph{Use the adjective lingering before effects.}
  \item \textbf{linger.} Several important questions continue to linger about the study's methodology. \emph{Use the base verb linger after continue to.}
\end{enumerate}

\subsubsection{Correct or incorrect?}
\begin{enumerate}
  \item \textbf{Correct} \emph{Linger in a place naturally describes a smell that remains noticeable.}
  \item \textbf{Incorrect}. The students lingered in the classroom after the lesson. \emph{Linger is normally intransitive, so use in before the place.}
  \item \textbf{Correct} \emph{Continue to linger is a common pattern for doubts or concerns that remain unresolved.}
  \item \textbf{Incorrect}. The economic effects lingered for several years. \emph{Use for to state how long something continues.}
  \item \textbf{Correct} \emph{Linger on naturally means keep attention focused on something for extra time.}
\end{enumerate}

\subsubsection{Collocation practice}
\begin{enumerate}
  \item \textbf{linger.} The effects of the recession continued to linger for several years. \emph{Effects can linger when they continue after the original event.}
  \item \textbf{after.} A few students lingered after class to speak with the professor. \emph{Linger after an event means remain after it has ended.}
  \item \textbf{linger.} Important doubts still linger about the reliability of the results. \emph{Doubts linger means uncertainty remains.}
  \item \textbf{over.} The author lingers over the social consequences of the policy. \emph{Linger over means spend extra time discussing or considering something.}
  \item \textbf{lingering.} The report examines the lingering effects of prolonged unemployment. \emph{Lingering effects are consequences that continue after the original cause or event.}
\end{enumerate}

\vfill
\begin{center}
\small\color{Muted} Created and compiled by \textbf{Amir Kooshky}\\
\href{https://t.me/KooshkyTOEFL}{Telegram Channel}\quad\textbullet\quad
\href{https://www.instagram.com/kooshkytoefl}{Instagram}\quad\textbullet\quad
\href{https://t.me/KooshkyTOEFL\_pv}{Personal Telegram}
\end{center}
\end{document}
